% ====================================================
% 부록 O – 집단 크기의 보편적 임계 스케일링
% ====================================================
\documentclass[12pt]{article}
\usepackage[utf8]{inputenc}
\usepackage[T1]{fontenc}
\usepackage[english]{babel}   % 영어 일부 사용

% ============================================================
% 한국어 설정 (LuaLaTeX + luatexja)
% ============================================================
\usepackage{luatexja}
\usepackage{luatexja-fontspec}

% 한국어 고딕체 (sans) – Noto Sans KR
\setsansjfont{Noto Sans KR}[
UprightFont = * Regular,
BoldFont = * Bold,
Script = Hangul,
Language = Korean
]

% 한국어 명조체 (serif) – Noto Serif KR
\IfFontExistsTF{Noto Serif KR}{
	\setmainjfont{Noto Serif KR}[
	UprightFont = * Regular,
	BoldFont = * Bold,
	Script = Hangul,
	Language = Korean
	]
}{
	% fallback: Noto Sans KR
	\setmainjfont{Noto Sans KR}[
	UprightFont = * Regular,
	BoldFont = * Bold,
	Script = Hangul,
	Language = Korean
	]
}

% 고정폭 폰트 – 라틴 문자는 Latin Modern Mono, 한글은 Noto Sans KR
\setmonojfont{Noto Sans KR}[
UprightFont = * Regular,
BoldFont = * Bold,
Script = Hangul,
Language = Korean
]

% 라틴 문자 폰트
\setmainfont{Times New Roman}[Ligatures=TeX]
\setsansfont{Arial}[Ligatures=TeX]
\setmonofont{Latin Modern Mono}[
WordSpace={1,0,0},
HyphenChar=None,
PunctuationSpace=WordSpace
]

% ============================================================
% 수학 및 서식 패키지
% ============================================================
\usepackage{amsmath,amssymb,amsthm,mathtools}
\usepackage{geometry}
\usepackage{booktabs}
\usepackage{hyperref}
\usepackage{url}
\usepackage{tabularx}
\usepackage{array}
\usepackage{makecell}
\usepackage{ragged2e}
\usepackage{bm}
\usepackage{graphicx}
\usepackage{caption}
\usepackage{subcaption}
\usepackage{float}

% 정리 환경 (한국어)
\newtheorem{theorem}{정리}
\newtheorem{lemma}{보조정리}
\newtheorem{definition}{정의}
\newtheorem{corollary}{따름정리}
\theoremstyle{remark}
\newtheorem{remark}{비고}

% 간단 명령어
\newcommand{\Keff}{K_{\mathrm{eff}}}
\newcommand{\eps}{\varepsilon}
\newcommand{\df}{d_{\!f}}
\newcommand{\kappac}{\kappa}
\newcommand{\betaf}{\beta}
\newcommand{\Tmin}{T_{\min}}
\newcommand{\Dprior}{D_{\mathrm{prior}}}

\hypersetup{colorlinks=true,linkcolor=blue,citecolor=blue,urlcolor=blue}

\begin{document}
	
	\begin{titlepage}
		\begin{center}
			\vspace*{0.5cm}
			\Huge\textbf{부록 O. 임계 집단 크기의 보편적 스케일링: \\ 경험적 증거와 기술적 임계 공리(YUCT)와의 연결}
			
			\vspace{0.3cm}
			\small\textit{생물학적 및 사회적 시스템에서 임계 집단 크기에 대한 체계적인 데이터 수집, 기술적 사전 임계 공리(부록 O)와 보편 지수 \(\beta = 0.67\)(부록 L)을 통해 해석됨.}
			
			\vspace{0.3cm}
			\small\textbf{알렉세이 V. 야쿠셰프} \\
			YUCT 이론: \url{https://yuct.org/}\\
			YPSDC 프로토콜: \url{https://ypsdc.com/}\\
			
			\vspace{2cm}
			\textbf{DOI:} \url{https://doi.org/10.5281/zenodo.18444598}\\
			
			\vspace{1cm}
			\large 
			이 작업의 단축 버전은 이전에 출판되었으며 다음에서 확인 가능\\
			\url{https://doi.org/10.5281/zenodo.18362308}\\
			\vspace{1cm}
			2026년 3월 \\ \large 버전 2.0.
			
			\vspace{0.5cm}
			\begin{minipage}{0.95\textwidth}
				\begin{abstract}
					\noindent
					이 부록은 다양한 집단 시스템이 군집, 분열 또는 구조 재편성과 같은 상전이를 겪는 임계 크기에 대한 경험적 데이터를 수집합니다. 조사된 시스템에는 꿀벌 군집, 군사 부대, 물고기 떼, 돌고래 무리, 가축 개체군 및 영장류 집단이 포함됩니다. 보편 오차 스케일링 법칙 \(\varepsilon \propto N^{\beta}\) (\(\beta = 0.67\), 부록 L)과 기술적 사전 임계 공리(부록 O)를 사용하여, 전이가 발생하는 임계 크기 \(N_{\text{crit}}\)와 최적 집단 크기 \(N_{\text{opt}}\) 사이의 정량적 관계를 유도합니다. 다양한 시스템에 대해 관찰된 비율 \(N_{\text{crit}}/N_{\text{opt}}\)은 \(\beta = 0.67\)과 일치하는 값 주변에 모여 있으며, 이는 이 지수의 보편성과 기술적(또는 조직적) 임계를 넘는 것이 조정 효율성의 상전이에 해당한다는 아이디어에 대한 독립적 지지를 제공합니다. 또한 YUCT 프레임워크를 더욱 검증할 수 있는 미래 실험 및 현장 연구를 위한 구체적이고 테스트 가능한 예측을 제안합니다.
				\end{abstract}
			\end{minipage}
			
			\vspace{0.3cm}
			\noindent
			\small\textbf{키워드:} YUCT, 임계 집단 크기, 상전이, 꿀벌 군집, 군사 계층, 물고기 떼, 보편 스케일링, \(\beta = 0.67\), 기술적 임계.
		\end{center}
	\end{titlepage}
	
	\tableofcontents
	\newpage
	
	\section{서론}
	\label{sec:intro}
	
	야쿠셰프 통합 조정 이론(YUCT)의 핵심 예측 중 하나는 시스템 크기에 따른 조정 오차의 스케일링을 지배하는 보편 지수 \(\beta = 0.67\)의 존재입니다(부록 L). 이 지수는 사전의 프랙탈 구조와 정보 이론적 제약에서 비롯되며, DNA 복제에서 우주 마이크로파 배경 변동에 이르기까지 40자릿수에 걸쳐 검증되었습니다. \(\beta\)가 진정으로 보편적이라면, 군집, 분열 또는 계층적 재편성과 같은 상전이를 겪는 집단 시스템의 임계 크기에서도 나타나야 합니다. 더욱이 기술적 사전 임계 공리(부록 O)는 시스템이 사전을 사용하거나 생성하기 위해 최소 기술 수준 \(\Tmin\)에 도달해야 한다고 명시하며; 이 임계는 시스템의 조정 효율성 \(\Keff\)의 임계점으로 재해석될 수 있습니다. 이 부록은 다양한 분야의 경험적 데이터를 수집하고 관찰된 임계 대 최적 집단 크기 비율이 YUCT의 예측과 일치함을 보여줌으로써 추상적 공리를 구체적이고 측정 가능한 양과 연결합니다.
	
	\section{이론적 틀}
	\label{sec:theory}
	
	크기 \(N\)의 조정된 시스템에 대해 상대 조정 오차 \(\varepsilon\)는 다음과 같이 스케일링된다고 가정합니다.
	\begin{equation}
		\varepsilon(N) = \varepsilon_0 N^{\beta}, \qquad \beta = 0.67,
		\label{eq:scaling}
	\end{equation}
	여기서 \(\varepsilon_0\)는 시스템 의존 상수입니다. 시스템은 \(\varepsilon\)가 최소인 크기 \(N_{\text{opt}}\)에서 최적으로 작동합니다. \(N\)이 \(N_{\text{opt}}\)를 초과함에 따라 오차가 증가하고, 임계 임계값 \(\varepsilon_{\text{crit}}\)에 도달하면 시스템은 상전이를 겪습니다(예: 하위 그룹으로 분열). 따라서
	\begin{equation}
		\varepsilon(N_{\text{crit}}) = \varepsilon_{\text{crit}}.
	\end{equation}
	전이 후 딸 그룹의 크기가 최적에 가깝다면, 수 보존에 의해 \(N_{\text{crit}} \approx k N_{\text{opt}}\) (\(k > 1\))입니다. \(k\)의 정확한 값은 전이의 성격(대칭 또는 비대칭 분열)에 따라 달라집니다. (\ref{eq:scaling})을 사용하면 다음을 얻습니다.
	\begin{equation}
		k^{\beta} = \frac{\varepsilon_{\text{crit}}}{\varepsilon_{\text{opt}}}, 
		\qquad\text{따라서}\qquad 
		k = \left( \frac{\varepsilon_{\text{crit}}}{\varepsilon_{\text{opt}}} \right)^{1/\beta}.
		\label{eq:k}
	\end{equation}
	따라서 비율 \(k\)는 임계 오차 대 최적 오차의 비율을 \(1/\beta\) 제곱한 것에 의해서만 결정됩니다. \(\varepsilon_{\text{crit}}/\varepsilon_{\text{opt}}\)가 다양한 시스템에서 대략 일정하다면(기본적인 물리적 또는 생물학적 제약에 의해 임계가 설정된다는 그럴듯한 가설), \(k\)도 보편적이어야 하며, 그 값은 경험적 데이터와 비교될 수 있습니다.
	
	\section{임계 집단 크기에 대한 경험적 데이터}
	\label{sec:data}
	
	우리는 집단 크기와 전이점이 문서화된 여러 잘 연구된 시스템에서 데이터를 수집했습니다.
	
	\subsection{꿀벌 군집 (Apis mellifera)}
	\label{subsec:bees}
	
	양봉 관행은 신뢰할 수 있는 수치를 제공합니다:
	\begin{itemize}
		\item 최적 군집 크기(첫 번째 벌떼, ``주 벌떼''): \(N_{\text{opt}} = 50\,000\)–\(60\,000\) 일벌(평균 질량 7–8 kg, 개체 질량 \(\approx 0.12\) g). \cite{Seeley2010, Winston1987}
		\item 벌떼 발생 전 군집 크기: \(N_{\text{crit}} = 80\,000\)–\(100\,000\) (때로는 더 많음). \cite{bees}
		\item 따라서 \(k_{\text{bees}} = N_{\text{crit}}/N_{\text{opt}} \approx 1.5\)–\(1.67\).
	\end{itemize}
	벌떼는 비대칭적입니다: 주 벌떼는 늙은 여왕벌을 데리고 더 작은 잔여 군집을 남깁니다. 관찰된 \(k\)는 약 1.6입니다.
	
	\subsection{군사 부대}
	\label{subsec:military}
	
	현대 군사 조직은 잘 정의된 부대 크기를 가진 명확한 계층 구조를 보입니다 \cite{FM3-90.5}:
	\begin{itemize}
		\item 분대: 8–12명
		\item 소대: 30–50명
		\item 중대: 100–200명
		\item 대대: 300–1000명
		\item 여단: 3000–5000명
		\item 사단: 9000–18000명
	\end{itemize}
	연속적인 수준 간 비율은 3에서 5까지 다양하며, 3–4 주변에 모여 있습니다. 중대의 경우, ~200을 초과하면 두 개의 소대로 분할될 수 있습니다. 이것은 \(k \approx 2\)를 제공하며, 이는 수준 간 비율보다 작습니다. 실제 분열 사건에 대한 더 나은 데이터가 필요하지만, 명확한 계층적 수준의 존재 자체는 조직적 복잡성의 양자화된 도약을 시사합니다.
	
	\subsection{물고기 떼 – 제브라피쉬 (Danio rerio)}
	\label{subsec:fish}
	
	제브라피쉬의 집단 행동에 대한 통제된 실험 \cite{Tunstrom2013}은 밀도가 증가함에 따라 극성화된 무리 짓기에서 무질서한 회전으로의 상전이를 보여줍니다. 임계 밀도는 표준 수조에서 약 50–100마리의 물고기 그룹 크기에 해당합니다. 정확한 \(N_{\text{opt}}\)는 제공되지 않지만, 전이는 그룹이 특정 크기를 초과할 때 발생합니다. 이는 \(k\)가 약 2라는 아이디어와 일치합니다.
	
	\subsection{돌고래 무리 (Tursiops spp.)}
	\label{subsec:dolphins}
	
	큰코돌고래에 대한 현장 연구 \cite{Connor2000}는 30–70개체의 전형적인 무리 크기를 보고합니다. 포식자(상어)가 있을 때, 무리는 수백 마리(최대 600)까지 모일 수 있습니다. 이 모임은 방어적 전이지 분열은 아니지만, 조정의 임계 변화를 보여줍니다. 큰 무리 대 전형적인 무리의 비율은 약 5–10이며, 이는 반대 방향(합체)의 \(k\)에 해당할 수 있습니다.
	
	\subsection{가축 개체군 – 유전적 보존}
	\label{subsec:livestock}
	
	유전적 보존을 위한 임계 개체군 크기에 대한 식량농업기구(FAO)의 데이터 \cite{FAO2013}:
	\begin{itemize}
		\item 소: 개체군 <2000일 때 임계 상태, 5000–15000에서 취약, >25000 정상.
		\item 양: 유전자 풀 보존을 위해 100–250 암양 권장.
		\item 돼지: 수컷 25 + 암컷 100.
		\item 닭: 수탉 50 + 암탉 250.
	\end{itemize}
	이 숫자는 최소 생존 가능 개체군 크기를 반영하며, 이는 유전적 조정(근친 교배 억제 회피)의 임계 임계값으로 볼 수 있습니다. \(N_{\text{opt}}\)를 최적 유전 건강을 위한 크기(예: 소의 경우 25000)로, \(N_{\text{crit}}\)를 최소 생존 가능(예: 2000)으로 해석하면 \(k \approx 0.08\)을 얻으며, 이는 1보다 작습니다 – 이것은 다른 종류의 전이(멸종)입니다. 그럼에도 불구하고 비율은 어떤 방식으로 \(\beta\)와 스케일링될 수 있지만, 별도의 모델이 필요합니다.
	
	\subsection{영장류 무리}
	\label{subsec:primates}
	
	비비 무리에 대한 데이터 \cite{Alberts1995}는 30–80개체의 전형적인 크기를 보여줍니다. 무리가 너무 커지면(약 100 이상), 두 개의 딸 무리로 분열될 수 있습니다. 이것은 \(k \approx 2\)–\(3\)을 제공합니다.
	
	\section{YUCT 예측과의 비교}
	\label{sec:comparison}
	
	식 (\ref{eq:k})에 따르면, 비율 \(k = N_{\text{crit}}/N_{\text{opt}}\)는 오차 비율 \(\varepsilon_{\text{crit}}/\varepsilon_{\text{opt}}\)에 의존합니다. 이 오차 비율이 보편적이라면, \(k\)는 종 전체에 걸쳐 일정해야 합니다. 표 \ref{tab:summary}에 요약된 데이터는 \(k\)가 실제로 약 1.5–2.5의 비교적 좁은 범위에 있음을 시사합니다.
	
	\begin{table}[htbp]
		\centering
		\caption{다양한 시스템에 대한 임계 크기 비율 요약}
		\label{tab:summary}
		\begin{tabular}{lccc}
			\toprule
			\textbf{시스템} & \(N_{\text{opt}}\) & \(N_{\text{crit}}\) & \(k = N_{\text{crit}}/N_{\text{opt}}\) \\
			\midrule
			꿀벌 군집 & 50–60k & 80–100k & 1.5–1.7 \\
			군사 부대 (소대→중대) & 30–50 & 100–200 & 2–4 \\
			제브라피쉬 떼 & 50? & 100? & 2? \\
			비비 무리 & 30–80 & ~100 & 1.3–3 \\
			\bottomrule
		\end{tabular}
	\end{table}
	
	대표 값으로 \(k \approx 1.6\)(벌, 영장류)을 취하면, 식 (\ref{eq:k})에서 다음을 얻습니다.
	\[
	\frac{\varepsilon_{\text{crit}}}{\varepsilon_{\text{opt}}} = k^{\beta} = 1.6^{0.67} \approx 1.37.
	\]
	이는 임계 오차가 최적 오차보다 약 37% 더 크다는 것을 의미합니다 – 상전이를 시작하기에 그럴듯한 임계값입니다. \(k \approx 2\)(물고기, 일부 영장류 무리)에 대해 \(2^{0.67} \approx 1.59\)를 얻습니다. 이 값들의 일관성은 \(\varepsilon_{\text{crit}}/\varepsilon_{\text{opt}}\)가 실제로 시스템 전반에 걸쳐 대략 일정함을 시사하며, 이는 \(\beta\)의 보편성과 기술적/조직적 임계를 조정 효율성의 상전이로 해석하는 것을 지지합니다.
	
	\section{테스트 가능한 예측}
	\label{sec:predictions}
	
	위의 틀을 바탕으로 실험적 또는 관찰적으로 테스트할 수 있는 몇 가지 구체적인 예측을 공식화합니다.
	
	\begin{enumerate}
		\item \textbf{꿀벌:} 주 벌떼 크기 대 벌떼 발생 전 군집 크기의 비율은 다양한 아종 및 환경 조건에서 일정해야 하며, \(\beta = 0.67\)인 \(k = (\varepsilon_{\text{crit}}/\varepsilon_{\text{opt}})^{1/\beta}\)를 충족해야 합니다. \(\varepsilon_{\text{crit}}/\varepsilon_{\text{opt}}\)를 독립적으로 추정할 수 있다면(예: 채집 효율성 또는 질병 발생률로부터), 예측된 \(k\)는 관찰과 일치해야 합니다.
		
		\item \textbf{물고기 떼:} 제브라피쉬 또는 다른 종에 대한 통제된 실험에서, 무리 짓기에서 회전(또는 분열)으로의 전이가 발생하는 집단 크기는 동일한 법칙에 따라 최적 집단 크기와 스케일링되어야 합니다. 조정 오차(예: 방향의 분산) 측정은 \(\varepsilon_{\text{opt}}\) 및 \(\varepsilon_{\text{crit}}\)의 독립적 추정치를 제공할 수 있습니다.
		
		\item \textbf{영장류:} 비비, 마카크 또는 침팬지 무리에 대한 장기 현장 연구는 분열 사건과 부모 및 딸 그룹의 크기를 기록해야 합니다. 비율 \(k\)는 안정적이어야 하며 YUCT 예측과 일치해야 합니다.
		
		\item \textbf{군사 조직:} 군사 부대의 분할(예: 대대가 두 개로 분할)에 대한 역사적 데이터를 조사할 수 있습니다. 이러한 분할이 발생하는 크기는 동일한 스케일링을 따라야 하며, 계층적 특성으로 인해 다른 \(k\)를 가질 수 있습니다.
		
		\item \textbf{가축 개체군:} 유전적 보존을 위해, 유전적 다양성을 유지하는 데 필요한 유효 개체군 크기 \(N_e\)는 멸종을 위한 임계 크기와 관련될 수 있습니다. 최소 생존 가능 \(N_e\) 대 장기 생존을 위한 최적 \(N_e\)의 비율도 \(\beta\)와 스케일링될 수 있지만, 이는 더 추측적입니다.
		
		\item \textbf{소셜 네트워크:} 온라인 소셜 그룹(예: WhatsApp 채팅, Reddit 커뮤니티)은 특정 크기를 초과할 때 자주 분할됩니다. 분할을 위한 임계 크기와 딸 그룹의 전형적인 크기는 디지털 흔적을 사용하여 분석될 수 있으며, 이는 \(k\)와 \(\beta\)의 보편성을 테스트하기 위한 방대한 데이터 세트를 제공합니다.
	\end{enumerate}
	
	\section{기술적 임계 공리와의 연결 (부록 O)}
	\label{sec:connection}
	
	기술적 사전 임계 공리(부록 O)는 시스템이 사전을 사용하거나 생성하기 위해 최소 기술 수준 \(\Tmin\)에 도달해야 한다고 명시합니다. 집단 크기 전이의 맥락에서, \(\Tmin\)은 일관성을 유지하는 데 필요한 임계 조정 효율성 \(K_{\mathrm{eff},\text{crit}}\)로 재해석될 수 있습니다. 경험적 비율 \(k = N_{\text{crit}}/N_{\text{opt}}\)는 시스템이 분열 또는 재조직화가 필요하기 전에 얼마나 성장할 수 있는지의 척도가 됩니다. \(k\)가 다양한 시스템에서 약 1.5–2 정도라는 사실은 기본 \(K_{\mathrm{eff},\text{crit}}/K_{\mathrm{eff},\text{opt}}\) 비율이 프랙탈 스케일링 법칙과 일치하게 보편적임을 시사합니다. 따라서 여기에 제시된 데이터는 기술적 임계의 존재와 보편 지수 \(\beta\)와의 연결에 대한 간접적 지지를 제공합니다.
	
	\section{결론}
	\label{sec:conclusion}
	
	이 부록에 수집된 데이터는 비록 예비적이지만, 집단 시스템이 상전이를 겪는 임계 집단 크기가 임의적이지 않고 YUCT 지수 \(\beta = 0.67\) 및 기술적 사전 임계 공리와 연결된 보편적 스케일링 법칙을 따른다는 것을 시사합니다. 비율 \(k = N_{\text{crit}}/N_{\text{opt}}\)는 약 1.5–2.5 주변에 모여 있으며, 암시된 오차 임계 \(\varepsilon_{\text{crit}}/\varepsilon_{\text{opt}} \approx 1.3\)–1.6은 다양한 시스템에서 놀랍도록 일관됩니다. 우리는 미래 실험 및 현장 연구를 통해 검증될 수 있는 구체적이고 테스트 가능한 예측을 제시했습니다. 이러한 예측의 확인은 YUCT와 조정 효율성 메트릭 \(\Keff\) 및 그 스케일링 법칙의 근본적 성격에 대한 강력한 독립적 지지를 제공할 것입니다.
	
	\begin{thebibliography}{99}
		\bibitem{Seeley2010} Seeley, T. D. (2010). \emph{Honeybee Democracy}. Princeton University Press.
		\bibitem{Winston1987} Winston, M. L. (1987). \emph{The Biology of the Honey Bee}. Harvard University Press.
		\bibitem{FM3-90.5} FM 3-90.5 (2008). \emph{Combined Arms Battalion}. US Army Field Manual.
		\bibitem{Tunstrom2013} Tunstrøm, K. et al. (2013). Collective states, multistability and transitional behavior in schooling fish. \emph{PLoS Computational Biology}, 9(2), e1002915.
		\bibitem{Connor2000} Connor, R. C. et al. (2000). The bottlenose dolphin: social relationships in a fission‑fusion society. In: \emph{Cetacean Societies}, University of Chicago Press.
		\bibitem{FAO2013} FAO (2013). \emph{In vivo conservation of animal genetic resources}. FAO Animal Production and Health Guidelines No. 14.
		\bibitem{Alberts1995} Alberts, S. C. \& Altmann, J. (1995). Balancing costs and opportunities: dispersal in male baboons. \emph{American Naturalist}, 145(2), 279–306.
		\bibitem{Yakushev2026} Yakushev, A. V. (2026). \emph{Yakushev Unified Coordination Theory (YUCT)}. Zenodo. \url{https://doi.org/10.5281/zenodo.18444599}
		\bibitem{AppendixL} 부록 L: 프랙탈 조정 오차 스케일링 – DNA에서 우주론까지의 보편 법칙.
		\bibitem{AppendixO} 부록 O: 기술적 사전 임계 공리.
		\bibitem{bees} Seeley, T. D. (2010) – 위와 동일, 또는 특정 참고문헌으로 대체 가능.
	\end{thebibliography}
	
\end{document}